% !TEX root = ../thuthesis-main.tex

\chapter{引言}

\section{研究背景}

随着集成电路工艺节点的持续演进与片上系统（System on Chip, SoC）规模的急剧扩大，数字芯片的设计复杂度正以指数级增长。
在这一背景下，电子设计自动化（Electronic Design Automation, EDA）工具承担着将人类设计意图转化为可制造芯片的使命。
从逻辑综合到时序分析，从布局布线到物理验证，EDA工具链已在设计流程的各环节实现了高度自动化。
然而，一个长期被忽视的瓶颈在于设计流程的起点——将人类可读的视觉设计规范转化为机器可执行的硬件描述语言（Hardware Description Language, HDL）代码。这一步骤至今仍以手工完成为主。

数字电路时序图（Digital Timing Diagram）是这类视觉化设计规范中最具代表性的一种。
它以时钟周期为横轴，以信号电平为纵轴，清晰地编码了数字系统的行为规约：信号之间的先后次序、组合逻辑与时序逻辑的因果关系，以及不同功能状态之间的转换条件。
无论是数据手册中的通信协议时序（如SPI、I2C、UART），还是复杂总线协议（如AMBA AXI\inlinecite{arm2011axi}、AHB\inlinecite{arm2021ahb}、APB\inlinecite{arm2021apb}）的传输波形，时序图都是定义模块行为的权威依据。
可以说，时序图本身就是一种“图形化的硬件规格说明书”，其中蕴含了实现对应HDL代码所需的全部逻辑信息。

当前的工程实践中，将时序图转化为HDL代码的过程完全依赖工程师的个人经验与手工劳动。
工程师需要逐周期地阅读波形，在头脑中重建有限状态机（Finite State Machine, FSM）与数据通路，再将这些抽象逻辑翻译为Verilog或VHDL代码。
这一过程存在两个根本性的问题：
其一，效率低下——包含数十个信号、跨越数百个时钟周期的复杂协议时序图，人工解读往往需要数小时甚至更长时间；
其二，容易出错——信号间的细微依赖关系与边界条件极易被忽略，由此引发的设计缺陷通常要到仿真甚至流片阶段才被暴露，造成高昂的修复成本。
因此，如果能够自动地从时序图图像中提取结构化的时序信息，并据此生成功能正确的HDL代码，将从根本上改变芯片设计的前端范式。

近年来，深度学习技术的飞速发展为这一问题带来了全新的解决思路。
大语言模型（Large Language Model, LLM）以Transformer架构\inlinecite{vaswani2017attention}为基础，经历了从GPT系列\inlinecite{radford2019language,brown2020language,achiam2023gpt4}到LLaMA\inlinecite{touvron2023llama}、Qwen3\inlinecite{yang2025qwen3}等开源模型的快速迭代，在自然语言理解、代码生成等任务上已展现出接近甚至超越人类的表现。
在硬件设计领域，已有多项工作利用LLM实现从文本需求到Verilog代码的自动生成\inlinecite{liu2023verilogeval,thakur2023benchmarking,thakur2024verigen,liu2023chipnemo,fu2023gpt4aigchip,blocklove2023chip,llm4hwdesign}，证明了LLM在理解硬件语义方面的潜力。
与此同时，多模态大语言模型（Multimodal Large Language Model, MLLM）的兴起进一步打破了模态壁垒。
GPT-4o\inlinecite{openai2024gpt4o}、Gemini\inlinecite{team2023gemini}等闭源模型，以及LLaVA\inlinecite{liu2023llava}、InternVL\inlinecite{chen2024internvl}、Qwen2.5-VL\inlinecite{bai2025qwen25vl}、Qwen3-VL\inlinecite{bai2025qwen3vl}等开源模型，已在自然图像理解、文档解析、图表推理等任务中取得了卓越成效。
这些进展表明，MLLM具备从视觉输入中提取结构化语义信息的基础能力，为实现“时序图图像到HDL代码”的自动化提供了技术基础。

然而，将通用MLLM直接应用于数字电路时序图的解析效果并不理想。
在时序图的自动化处理方面，已有若干相关研究。
He等人\inlinecite{he2023tdmagic}提出的TD-Magic能够从时序图图像中提取严格偏序关系作为形式化规范，但其基于规则的多模块设计难以扩展到复杂场景，且仅适用于包含显式偏序关系的少数时序图。
同一团队后续提出的TD-Interpreter\inlinecite{he2025tdinterpreter}则采用了不同的技术路线，通过对开源MLLM LLaVA进行微调，构建了一个面向时序图的视觉问答（Visual Question Answering, VQA）系统。
TD-Interpreter开发了包含描述型与推理型两类数据的合成生成流程，使微调后的LLaVA在时序图理解任务上大幅超越了未经微调的GPT-4o。
该工作证明了通过领域特定数据对MLLM进行微调以理解时序图的可行性，但其定位为交互式问答工具——输入为时序图图像与自然语言问题，输出为自然语言回答，并不涉及结构化信息的提取与HDL代码的生成。
此外，Chang等人\inlinecite{chang2024natural}将MLLM引入硬件设计流程，利用视觉与文本的多模态输入生成Verilog代码，但该工作并未关注时序图，其视觉输入主要面向其他类型的设计图表。

与上述工作不同，本文选择将时序图解析任务建模为从图像到结构化表示、再到HDL代码的生成式流水线，而非视觉问答范式。这一任务设定的选择基于以下两方面的考量。

其一，在数据构建的可扩展性方面，结构化表示与HDL代码均可从现有的Verilog源代码通过仿真工具链自动获取——给定一个Verilog模块及其测试平台，仿真即可生成时序图图像，而模块源代码与仿真结果则天然构成对应的标注数据。这意味着整个数据构建流程可以完全自动化，无需像视觉问答任务那样依赖大量人工设计的问答对，从而具备更强的规模化能力。

其二，在任务本身的深度与实用性方面，结构化表示要求模型忠实地还原时序图中每一个信号的名称、电平序列与时钟关系，其输出可被程序化地解析与比较，从而为自动化评估提供客观基础，也为从时序图自动生成断言、检测设计文档与RTL实现的一致性等潜在下游应用奠定技术基础。在此基础上进一步生成HDL代码，本质上是对时序图所描述的硬件行为进行逆向工程——模型不仅需要”看懂”每个信号的波形，还需推理出信号之间的因果逻辑并将其转化为可综合的电路描述。因此，HDL代码的功能正确性也构成了衡量模型对时序图理解深度的严格标准。

尽管上述任务设定在数据构建的可扩展性与评价深度方面均具有优势，但要在此框架下实现端到端的自动化解析，仍需克服一系列关键技术瓶颈。具体而言，本文所面对的核心挑战包括：
第一，训练数据方面，当前不存在面向时序图解析的高质量标注数据集，且基于随机激励生成的时序图往往缺少有意义的功能模式，无法为模型训练提供有效监督信号；
第二，模型能力方面，时序图的视觉特征高度领域化，需要通过针对性的微调使MLLM掌握从像素到结构化时序信息的映射能力，同时还需训练文本LLM完成从结构化表示到可综合HDL代码的转换；
第三，评价体系方面，传统的文本相似度指标（如BLEU、ROUGE）无法反映结构化时序表示的语义正确性，亟需设计专门的评价方法；
第四，代码质量方面，仅依赖有监督学习难以保证生成代码的功能正确性，需要引入基于仿真反馈的优化机制来闭合“生成-验证”的循环。


\section{本文贡献}

针对上述挑战，本文提出了一种基于多模态大语言模型的数字电路时序图自动化解析框架。
该框架将时序图解析建模为两阶段任务：第一阶段由MLLM将时序图图像转换为结构化的中间表示，第二阶段由文本LLM将结构化表示翻译为对应的HDL代码。本文的主要贡献如下：

\begin{enumerate}
  \item \textbf{优化数据合成链路，构建高质量训练数据集。}
  针对该领域标注数据匮乏的问题，本文对数据合成链路进行了系统优化。
  在数据生成端，利用大语言模型替代随机激励来生成测试平台，使仿真产生的时序图呈现出更具工程意义的信号行为模式。
  在数据筛选端，设计了严格的质量过滤机制，剔除信号行为退化、缺乏区分度的低质量样本。
  在数据划分端，采用基于并查集的防泄漏划分策略，通过多维度相似性检测将高度相似的样本归入同一等价类并整组分配，避免训练集与测试集之间的信息泄漏，确保评估结果的可靠性。

  \item \textbf{提出面向时序图结构化表示的专用评价指标。}
  鉴于BLEU、ROUGE等通用文本相似度指标对信号排列顺序敏感且无法区分语义关键错误的严重程度，本文设计了基于信号级精确率-召回率的专用评价指标。
  该指标首先通过精确匹配与模糊匹配建立预测信号与标准答案信号之间的对应关系，然后在展开保持符号后逐周期比较波形序列的一致性，最终汇总为F1分数。该指标同时作为强化学习阶段的奖励函数，直接驱动模型优化功能正确性。

  \item \textbf{通过两阶段有监督微调建立基线能力。}
  本文采用有监督微调（SFT）分别训练两个阶段的模型：
  对多模态大语言模型进行微调，使其从时序图图像中提取信号名称、电平序列及时钟关系并输出为结构化表示；
  对纯文本大语言模型进行微调，使其将结构化表示转换为语法正确的HDL代码。
  两阶段的解耦设计降低了单一模型的学习难度，并使每个阶段都能独立优化。

  \item \textbf{通过基于仿真奖励的强化学习提升代码功能正确性。}
  SFT的词元级损失函数不直接优化生成代码的功能行为。本文将所提出的评价指标作为奖励函数，采用GRPO算法\inlinecite{shao2024deepseekmath}，以仿真对拍的F1分数作为连续值奖励信号对模型进行强化学习优化，形成”生成-仿真-评分-更新”的闭环，在SFT基线之上进一步提升代码的功能正确性。
\end{enumerate}

实验结果表明，本文所提出的方法在数字电路时序图解析任务上显著优于现有基线方法，验证了从时序图图像到结构化表示、再到HDL代码这一自动化流水线的可行性与有效性，为智能EDA工具的发展开辟了新的路径。


\section{论文组织结构}

本文其余部分的组织结构如下：

第二章介绍问题背景，包括大语言模型与多模态大语言模型的架构原理、数字电路时序图的WaveJSON表示格式、评价指标与有监督微调及强化学习的核心方法、面向硬件设计与时序图分析的已有工作，以及本文的问题定义。

第三章详细阐述本文提出的时序图自动化解析框架，涵盖数据合成链路的优化设计、多层级质量过滤与防泄漏数据划分、专用评价指标的构建方法、两阶段有监督微调策略以及基于强化学习的进一步优化。

第四章展示实验结果与分析，包括数据集统计分析、评价指标有效性验证、两阶段模型的定量评估与生成代码的错误模式分析等。

第五章总结全文工作，讨论现有方法的局限性，并展望未来的研究方向。
